% New proof increment. Do not automatically merge into the checkpoint.
% Unique label prefix: inc:AN02:arrival:v1:
\documentclass[11pt]{article}
\usepackage[T1]{fontenc}
\usepackage[utf8]{inputenc}
\usepackage{lmodern}
\usepackage[margin=0.95in]{geometry}
\usepackage{amsmath,amssymb,amsthm,mathtools,microtype}
\usepackage{booktabs,array,xurl}
\usepackage[hidelinks]{hyperref}
\usepackage{fancyhdr}
\pagestyle{fancy}
\fancyhf{}
\fancyhead[L]{\small AN02: initialized weak-side arrival}
\fancyhead[R]{\small Version 1 --- review increment}
\fancyfoot[C]{\thepage}
\setlength{\headheight}{14pt}
\setlength{\emergencystretch}{3em}
\allowdisplaybreaks[2]
\numberwithin{equation}{section}
\newtheorem{theorem}{Theorem}[section]
\newtheorem{lemma}[theorem]{Lemma}
\newtheorem{proposition}[theorem]{Proposition}
\newtheorem{corollary}[theorem]{Corollary}
\theoremstyle{remark}
\newtheorem{remark}[theorem]{Remark}
\newcommand{\E}{\mathbb E}
\newcommand{\R}{\mathbb R}
\newcommand{\dd}{\,d}
\newcommand{\ind}{\mathbf1}
\newcommand{\norm}[1]{\left\lVert#1\right\rVert}
\DeclareMathOperator{\tr}{tr}
\title{An initialized weak-side reservoir\\at a deterministic early envelope time}
\author{Separate analytical proof increment for AN2}
\date{October 2, 2026 --- version 1, for independent review}
\begin{document}
\maketitle

\begin{abstract}
We construct an initial-label interval, specified analytically from the
parameters rather than selected from a trained trajectory, whose exact
finite-noise population characteristics lie in an incoming weak-side sector
at a logarithmic early time. The original-flow mass of this interval has an
explicit lower bound. In particular, for
$3/5\le\rho\le3/4$, $0<q\le1/20$, and
$0<S_0<\delta\le1/100$, at
$T_\delta=(1+2q^2)^{-1}\log(\delta/S_0)$ it is at least
$(q/3000)S_0(S_0/\delta)^{1/16}$.
The construction uses a one-sided Gaussian-tail bound on backward angular
compression and the checkpoint's exact early comparison. It retains all
population interactions in transferring to the true flow.
The cohort is not yet captured in a bounded weak noise-scaled chart, and
$T_\delta$ is not a strong-learning hitting time. Thus this is a new
initialized arrival sublemma, not completion of AN2, AN3, or Theorems A/B.
\end{abstract}

\section{Contract, setting, and main result}
\label{inc:AN02:arrival:v1:contract}

\paragraph{Dependencies and status.}
The input model is exactly the aligned-balanced isotropic population in
\cite{F}, P1--P5, and \cite{P}, \texttt{pa:setup}.
The early comparison used below is \cite{P},
\texttt{pa:earlybounds} and \texttt{pa:targetcomparison}; the part needed
here is restated and verified in Section~\ref{inc:AN02:arrival:v1:early}.
The Gaussian feature identity comes from \cite{F}, P4, equation (6.7).
No leading-system approximation, coherent reduction, numerical orbit, or
assumption of a pre-existing weak attractor is used. Global existence of
the exact labeled characteristic solution is imported from \cite{F}, P5.
All new results are for review; the protected source documents are unchanged.

Use normalized time $\tau=\mu_1^2t_{\rm phys}$ and
\[
 P_1=\mathcal N(e_1,q^2I_2),\qquad
 P_2=\mathcal N(\rho e_2,q^2I_2),\qquad P=(P_1+P_2)/2.
\]
Let $\lambda=\rho^2$, $a_*=1/2+q^2$, and
$\lambda_0(d\alpha)=d\alpha/(2\pi)$. With
$a_\theta=(\cos\theta,\sin\theta)$ and
$b_\psi=(\cos\psi,\sin\psi)$, the exact population is
\[
 f_\tau(X)=\int m_\tau(\alpha)b_{\psi_\tau(\alpha)}
       (a_{\theta_\tau(\alpha)}^\top X)_+\dd\lambda_0(\alpha),
 \qquad \theta_0=\psi_0=\alpha,\quad m_0=S_0.
\]
Write $S(\tau)=\int m_\tau\dd\lambda_0$ and
\[
 R_f(\theta)=\E_P[(X-f(X))X^\top\ind_{\{a_\theta^\top X>0\}}].
\]
The full, untied, balanced equations are
\begin{equation}\label{inc:AN02:arrival:v1:flow}
 \theta'=b^\top R_fa^\perp,\qquad
 \psi'=(b^\perp)^\top R_fa,\qquad
 (\log m)'=2b^\top R_fa.
\end{equation}
Let $\phi,\Phi$ be the standard normal density and distribution function,
and set $K(r)=\phi(r)+r\Phi(r)=\E(Z+r)_+$.

\paragraph{An analytical scalar comparison, not a fitted trajectory.}
Put
\[
 A_\theta=\E_P[XX^\top\ind_{\{a_\theta^\top X>0\}}],\qquad
 v_0(\theta)=(a_\theta^\perp)^\top A_\theta a_\theta.
\]
Denote the scalar flow of $\widehat\theta'=v_0(\widehat\theta)$ by
$\varphi_\tau(\alpha)$. At positive $q$ this is a smooth flow on the circle;
we use its real lift. Fix once and for all the terminal sector
\begin{equation}\label{inc:AN02:arrival:v1:sector}
 I=[2\pi/3,3\pi/4].
\end{equation}
For a prescribed $T\ge0$, the initial-label interval is
\begin{equation}\label{inc:AN02:arrival:v1:cohort}
 C_T=\{\varphi_{-T}(\theta):\theta\in I\}.
\end{equation}
This interval is specified from $\rho,q,T$ before following the nonlinear
population. For each fixed $T$ it is held fixed on the whole interval
$[0,T]$. No empirical specialization labels enter its definition.

\begin{theorem}[Initialized incoming weak-side mass]
\label{inc:AN02:arrival:v1:main}
Assume
\begin{equation}\label{inc:AN02:arrival:v1:parameters}
 \frac35\le\rho\le\frac34,\qquad 0<q\le\frac1{20},\qquad
 0<S_0<\delta\le\frac1{100}.
\end{equation}
Choose any $R\ge1$ such that $qR\le\lambda/2$, and define
\begin{align}
 T_\delta&=\frac1{2a_*}\log\frac{\delta}{S_0},
 \label{inc:AN02:arrival:v1:clock}\\
 \kappa_R(q,\rho)&=\frac12\left[
 (1+2q^2)\Phi\!\left(-\frac1{2q}\right)
 +(\lambda+2q^2)\Phi(-R)\right],
 \qquad \varepsilon_R=\frac{\kappa_R}{2a_*}.
 \label{inc:AN02:arrival:v1:kappa}
\end{align}
For $C=C_{T_\delta}$, let
$M_{\rm in}=\int_C m_{T_\delta}(\alpha)\dd\lambda_0(\alpha)$.
Then the following are statements about the exact original flow:
\begin{align}
 S(\tau)&\le\delta\quad(0\le\tau\le T_\delta),
 \label{inc:AN02:arrival:v1:smallmass}\\
 \theta_{T_\delta}(\alpha),\ \psi_{T_\delta}(\alpha)
 &\in[7\pi/12,5\pi/6]\quad(\alpha\in C),
 \label{inc:AN02:arrival:v1:angles}\\
 m_{T_\delta}(\alpha)&\ge e^{-4\delta}S_0\quad(\alpha\in C),
 \label{inc:AN02:arrival:v1:labelmass}\\
 M_{\rm in}&\ge
 \frac{\rho qK(-R)}{72a_*}e^{-4\delta}
 S_0\left(\frac{S_0}{\delta}\right)^{\varepsilon_R}.
 \label{inc:AN02:arrival:v1:sharpseed}
\end{align}
In particular $R=1$ is admissible throughout
\eqref{inc:AN02:arrival:v1:parameters}, and a completely elementary bound is
\begin{equation}\label{inc:AN02:arrival:v1:rationalseed}
 \boxed{\displaystyle
 M_{\rm in}\ge\frac q{3000}\,S_0
       \left(\frac{S_0}{\delta}\right)^{1/16}.}
\end{equation}
The cohort's input and output directions both satisfy
$a_1,b_1\le-1/4$ and $a_2,b_2\ge1/2$ at $T_\delta$.
\end{theorem}

\paragraph{What is not in this statement.}
There is no claim that $S(T_\delta)=\delta$, that either concept is
substantially learned, or that the full incoming mass remains available
at a later resident-entry time. In particular $M_{\rm in}$ is not yet
$N_{\rm ent}$ in AN3. Directions in
\eqref{inc:AN02:arrival:v1:angles} are outside a bounded weak noise-scaled
chart as $q\to0$. Capture, subsequent amplification, and the
strong-population shape estimates are separate open obligations.

\section{The exact early comparison used in the transfer}
\label{inc:AN02:arrival:v1:early}

\begin{lemma}[Imported comparison, in the required normalization]
\label{inc:AN02:arrival:v1:comparison}
Define
\[
 \overline S(\tau)=S_0e^{2a_*\tau},\qquad
 \epsilon_{\rm al}(\tau)=\overline S(\tau)-S_0,
 \qquad
 H(\tau)=S_0(2e^{2a_*\tau}-3e^{a_*\tau}+1).
\]
While $\epsilon_{\rm al}<\pi$, the exact initialized population and the
target-only characteristic satisfy
\begin{align}
 S&\le\overline S,\qquad |\psi-\theta|\le\epsilon_{\rm al},
 \label{inc:AN02:arrival:v1:comparison1}\\
 |\theta-\varphi_\tau(\alpha)|&\le H,\qquad
 |\psi-\varphi_\tau(\alpha)|\le H+\epsilon_{\rm al},\qquad
 \left|\log\frac{m}{\widehat m}\right|\le2H,
 \label{inc:AN02:arrival:v1:comparison2}
\end{align}
where $\widehat m_0=S_0$ and
$(\log\widehat m)'=2g_0(\varphi_\tau(\alpha))$ with
$g_0(\theta)=a_\theta^\top A_\theta a_\theta\ge0$.
Consequently $\widehat m\ge S_0$.
\end{lemma}
\begin{proof}
This is \cite{P}, \texttt{pa:earlybounds} and
\texttt{pa:targetcomparison}; the following verifies the dependencies used
here. Put
$C_f(\theta)=\E_P[f(X)X^\top\ind_{\{a_\theta^\top X>0\}}]$.
The spectral bound $0\preceq A_\theta\preceq a_*I$ and the feature
representation give $\norm{C_f}\le a_*S$ by Cauchy--Schwarz.
The exact mass identity yields
\[
 S'=2\E_P[X^\top f]-2\E_P|f|^2\le2a_*S,
\]
so $S\le\overline S$. If $d=\psi-\theta$, symmetry of $A_\theta$ gives
\[
 d'=-\tr(A_\theta)\sin d+\eta,\qquad |\eta|\le2a_*S.
\]
Before the first possible exit from $|d|<\pi$, its target term points
inward, so $D^+|d|\le2a_*\overline S$. Integration closes that first-exit
condition whenever $\epsilon_{\rm al}<\pi$.

The moving-gate derivative obeys $A_\theta'a_\theta=0$: its boundary
integrand contains $X^\top a_\theta=0$. Thus
\[
 v_0'=(a^\perp)^\top A_\theta a^\perp-a^\top A_\theta a,
 \quad g_0'=2v_0,\quad |v_0'|,|g_0'|\le a_*.
\]
The last bound for $g_0'$ follows by comparing the two quadratic forms
on $(a+a^\perp)/\sqrt2$ and $(a-a^\perp)/\sqrt2$.
The true input angular velocity differs from $v_0(\theta)$ by at most
$a_*(\epsilon_{\rm al}+\overline S)$. The scalar comparison for the
angular distance is therefore
\[
 D^+|\theta-\widehat\theta|
 \le a_*|\theta-\widehat\theta|
       +a_*(\epsilon_{\rm al}+\overline S).
\]
Its zero-initial-data equality solution is precisely $H$. Moreover
\[
 |b^\top R_fa-g_0(\widehat\theta)|
 \le a_*(\epsilon_{\rm al}+\overline S+H)=H'.
\]
Integration proves the logarithmic mass comparison. The output-angle
comparison follows by adding the alignment error. Finally $g_0\ge0$
proves $\widehat m\ge S_0$. All remainder estimates include the entire
population through $C_f$, not merely the chosen cohort.
\end{proof}

\section{A Gaussian backward-transport estimate}
\label{inc:AN02:arrival:v1:transport}

\subsection{The scalar field and its upstream interval}
The exact Gaussian feature formula gives
\begin{equation}\label{inc:AN02:arrival:v1:vformula}
 v_0(\theta)=\frac q2\left[
 -\sin\theta K\!\left(\frac{\cos\theta}{q}\right)
 +\rho\cos\theta K\!\left(\frac{\rho\sin\theta}{q}\right)
 \right].
\end{equation}
The isotropic $q^2a_\theta$ part of each Gaussian feature vector is
orthogonal to $a_\theta^\perp$ and contributes nothing to this formula.

\begin{lemma}[An upstream barrier and two quantitative field bounds]
\label{inc:AN02:arrival:v1:barrier}
Assume the $\rho,q,R$ hypotheses of Theorem~\ref{inc:AN02:arrival:v1:main}.
Put
\[
 \beta=\pi+\arcsin\frac{qR}{\rho},\qquad
 B=\frac{4\pi}{3},\qquad
 b_R=\frac{\rho qK(-R)}6.
\]
Then $\pi<\beta<B$,
\begin{align}
 -v_0(\theta)&\ge b_R>0
       &&(2\pi/3\le\theta\le\beta),
 \label{inc:AN02:arrival:v1:speed}\\
 v_0' (\theta)&\le\kappa_R
       &&(\beta\le\theta\le B),
 \label{inc:AN02:arrival:v1:tailcurvature}\\
 |v_0(\theta)|&\le a_*/2
       &&(\theta\in\R),
 \label{inc:AN02:arrival:v1:maxspeed}
\end{align}
and $v_0(B)>0$.
Consequently there is a first zero $\theta_*\in(\beta,B)$ after $\beta$,
with $v_0<0$ on $[2\pi/3,\theta_*)$.
No uniqueness or nondegeneracy of zeros elsewhere is asserted.
\end{lemma}
\begin{proof}
Since $qR/\rho\le\rho/2\le3/8<1/2$, the definition of $\beta$ is valid
and $\beta<\pi+\pi/6<B$. For
$2\pi/3\le\theta\le\pi$, put $c=\cos\theta$, $s=\sin\theta$.
Here $-c\ge1/2$, $s\ge0$, and \eqref{inc:AN02:arrival:v1:vformula} gives
\[
 -v_0(\theta)\ge\frac{\rho q(-c)}2 K(\rho s/q)
 \ge\frac{\rho q}4K(0)\ge b_R.
\]
For $\pi\le\theta\le\beta$, write $c=-C$, $s=-d$.
Then $0\le d\le\rho/2$, $C\ge\sqrt3/2$, and $C\ge\rho d$.
Since $K$ is increasing and $\rho d/q\le R$,
\begin{align*}
 -v_0(\theta)
 &=\frac q2[\rho C K(-\rho d/q)-dK(-C/q)]\\
 &\ge\frac q2(\rho C-d)K(-\rho d/q)
 \ge\frac q2\rho\frac{\sqrt3-1}{2}K(-R)
 \ge b_R.
\end{align*}
The final inequality uses $(\sqrt3-1)/2>1/3$.

For the derivative bound, consider a Gaussian $X=c_p+qZ$ and let
$\ell=a_\theta^\top c_p$. Resolve $X$ along $a_\theta$ and its
orthogonal direction. Direct integration of the one-dimensional
truncated second moment gives
\begin{equation}\label{inc:AN02:arrival:v1:trace}
 \E[|X|^2\ind_{\{a_\theta^\top X>0\}}]
 =(|c_p|^2+2q^2)\Phi(\ell/q)+\ell q\phi(\ell/q).
\end{equation}
On $[\beta,B]$ the two projected means satisfy
$\ell_1=\cos\theta\le-1/2$ and
$\ell_2=\rho\sin\theta\le-qR$.
The final term in \eqref{inc:AN02:arrival:v1:trace} is nonpositive.
Because $v_0'\le\tr(A_\theta)$ by the derivative identity in
Lemma~\ref{inc:AN02:arrival:v1:comparison}, averaging the two Gaussian
bounds gives \eqref{inc:AN02:arrival:v1:tailcurvature}.
This step uses an upper bound on a signed derivative; no absolute
$1/q$-Lipschitz estimate is introduced.

For \eqref{inc:AN02:arrival:v1:maxspeed}, the difference between the
quadratic forms of $A_\theta$ on
$(a+a^\perp)/\sqrt2$ and $(a-a^\perp)/\sqrt2$ is $2v_0$.
Both forms lie in $[0,a_*]$.
Finally
\[
 v_0(B)=\frac q4\left[
 \sqrt3 K\!\left(-\frac1{2q}\right)
 -\rho K\!\left(-\frac{\sqrt3\rho}{2q}\right)\right]>0.
\]
Indeed $\sqrt3\rho>1$ on the specified ratio range, $\rho<\sqrt3$,
and $K$ is strictly increasing and strictly positive.
Continuity, the strict negative value at $\beta$, and the positive value
at $B$ yield the claimed first zero and sign interval.
\end{proof}

\subsection{A lower bound on surviving initial-label length}
\begin{lemma}[Backward Jacobian and initial interval size]
\label{inc:AN02:arrival:v1:jacobian}
For every $T\ge0$ and $\theta\in I$, the backward scalar characteristic
$\chi(T,\theta)=\varphi_{-T}(\theta)$ lies in
$[2\pi/3,\theta_*)$ and
\begin{equation}\label{inc:AN02:arrival:v1:jaclower}
 \partial_\theta\chi(T,\theta)
 \ge\frac{\rho qK(-R)}{3a_*}e^{-\kappa_R T}.
\end{equation}
Consequently
\begin{equation}\label{inc:AN02:arrival:v1:labellength}
 \lambda_0(C_T)\ge
 \frac{\rho qK(-R)}{72a_*}e^{-\kappa_R T}.
\end{equation}
More generally, for a subinterval $J\subset I$, replace $1/24=|I|/(2\pi)$
in this calculation by $|J|/(2\pi)$.
\end{lemma}
\begin{proof}
Put $w=-v_0$. The backward equation is $\partial_T\chi=w(\chi)$,
with $w>0$ on $[2\pi/3,\theta_*)$. It is increasing in backward time
and cannot reach the equilibrium $\theta_*$ at a finite time: otherwise
uniqueness, applied also backward from the contact time, would identify
it with the constant equilibrium solution. Thus every denominator below
is nonzero on the finite interval where it is used.

Differentiating the scalar flow gives
\begin{equation}\label{inc:AN02:arrival:v1:jacexact}
 \partial_\theta\chi(T,\theta)
 =\exp\!\left(-\int_0^T v_0'(\chi(s,\theta))\dd s\right)
 =\frac{w(\chi(T,\theta))}{w(\theta)}.
\end{equation}
The ratio equality follows either by differentiating both sides in $T$
or by differentiating $\log w(\chi)$; initially both sides equal one.
If $\chi(T,\theta)\le\beta$, then
$w(\chi(T,\theta))\ge b_R$ by
\eqref{inc:AN02:arrival:v1:speed}. If it has crossed $\beta$, let $s_b$
be that unique backward crossing time. Thereafter
\[
 \frac d{ds}\log w(\chi(s,\theta))
 =-v_0'(\chi(s,\theta))\ge-\kappa_R.
\]
Therefore in both cases
\[
 w(\chi(T,\theta))\ge b_Re^{-\kappa_R T}.
\]
Since $0<w(\theta)\le a_*/2$, substitution into
\eqref{inc:AN02:arrival:v1:jacexact} proves
\eqref{inc:AN02:arrival:v1:jaclower}. The flow is strictly increasing in
its starting coordinate, so the length of $C_T$ is the integral of this
Jacobian over $I$. Since $|I|=\pi/12$, division by $2\pi$ proves
\eqref{inc:AN02:arrival:v1:labellength}.
\end{proof}

\begin{remark}[Why the one-sided upstream bound matters]
Using only $|v_0'|\le a_*$ in the exponential Jacobian formula would give
$\lambda_0(C_T)\ge(1/24)e^{-a_*T}$. At $T=T_\delta$ this loses a factor
$(S_0/\delta)^{1/2}$ before accounting for any later capture.
Lemma~\ref{inc:AN02:arrival:v1:jacobian} replaces that exponent by
$\varepsilon_R$, with a displayed $qK(-R)$ prefactor. The finite passage
through the faster angular region is handled by the exact speed ratio,
not an exponential bound over the entire delay.
\end{remark}

\section{Transfer to the initialized nonlinear population}
\label{inc:AN02:arrival:v1:transfer}

\begin{proof}[Proof of Theorem~\ref{inc:AN02:arrival:v1:main}, except for the rational simplification]
For $0\le\tau\le T_\delta$, $\overline S(\tau)\le\delta$ and
$\epsilon_{\rm al}(\tau)\le\delta<\pi$.
Lemma~\ref{inc:AN02:arrival:v1:comparison} therefore applies on the
whole interval. At its endpoint,
\[
 H(T_\delta)=2\delta-3\sqrt{S_0\delta}+S_0\le2\delta,
 \qquad H(T_\delta)+\epsilon_{\rm al}(T_\delta)\le3\delta<\pi/12.
\]
For each label in the predetermined $C_{T_\delta}$,
$\varphi_{T_\delta}(\alpha)\in[2\pi/3,3\pi/4]$.
The comparison places both true angles in
$[7\pi/12,5\pi/6]$. On this interval sine is at least $1/2$ and cosine
is at most $-\sin(\pi/12)<-1/4$. Since
$\widehat m(T_\delta,\alpha)\ge S_0$, the mass comparison gives
$m(T_\delta,\alpha)\ge e^{-4\delta}S_0$.
Integrating over the fixed initial-label interval, and applying
Lemma~\ref{inc:AN02:arrival:v1:jacobian}, yields
\[
 M_{\rm in}\ge e^{-4\delta}S_0
       \frac{\rho qK(-R)}{72a_*}e^{-\kappa_R T_\delta}.
\]
The identity
$e^{-\kappa_R T_\delta}=(S_0/\delta)^{\varepsilon_R}$ completes the
sharp estimate. No density or monotonicity of the true input-angle
projection was assumed; its possible folds do not affect the labelwise
comparison or the fixed-label integral.
\end{proof}

\begin{proposition}[Elementary uniform constants]
\label{inc:AN02:arrival:v1:rational}
Under \eqref{inc:AN02:arrival:v1:parameters}, $R=1$ satisfies the theorem's
hypotheses, $\varepsilon_1<1/16$, and
\eqref{inc:AN02:arrival:v1:rationalseed} follows without numerical
Gaussian evaluation.
\end{proposition}
\begin{proof}
First $q\le1/20<9/50\le\rho^2/2$. The elementary bounds
\begin{equation}\label{inc:AN02:arrival:v1:elementary}
 \Phi(-1)<\frac3{16},\qquad K(-1)>\frac3{128},\qquad
 \Phi\!\left(-\frac1{2q}\right)\le4q^2\le\frac1{100}
\end{equation}
have short analytic proofs. Since $\phi(0)>3/8$ (for example using
$\pi<22/7$) and $e^{-z^2/2}\ge1-z^2/2$ for $0\le z\le1$,
\[
 \Phi(-1)=\frac12-\phi(0)\int_0^1e^{-z^2/2}\dd z
 \le\frac12-\frac56\phi(0)<\frac3{16}.
\]
Also
\[
 K(-1)=\int_1^\infty(z-1)\phi(z)\dd z
 \ge\frac12\phi(2)
 >\frac12\frac38\frac18=\frac3{128}.
\]
Here $e^2<8$, for instance by summing the exponential series through
degree four (sum $7$) and bounding its tail by
$(4/15)/(1-1/3)=2/5$. The final Gaussian bound in
\eqref{inc:AN02:arrival:v1:elementary} is Markov's inequality applied
to $Z^2$, since $\E Z^2=1$.

Because $2a_*=1+2q^2$ and $\lambda+2q^2\le227/400$,
\[
 \varepsilon_1
 =\frac12\Phi\!\left(-\frac1{2q}\right)
 +\frac{\lambda+2q^2}{2(1+2q^2)}\Phi(-1)
 \le\frac1{200}+\frac{227}{800}\frac3{16}
 =\frac{149}{2560}<\frac1{16}.
\]
The coefficient in \eqref{inc:AN02:arrival:v1:sharpseed} satisfies
\[
 \frac{\rho}{72a_*}e^{-4\delta}K(-1)
 \ge\frac{80}{67}\frac1{72}\frac{24}{25}\frac3{128}
 =\frac1{2680}>\frac1{3000},
\]
using $a_*\le201/400$ and $e^{-4\delta}\ge1-4\delta\ge24/25$.
Finally $S_0/\delta<1$, so
$(S_0/\delta)^{\varepsilon_1}\ge(S_0/\delta)^{1/16}$.
\end{proof}

\section{Exact training and probe meaning of the incoming cohort}
\label{inc:AN02:arrival:v1:meaning}
These conclusions identify what has actually been supplied to the next
capture problem. They do not replace that problem by a mass label.

At $T=T_\delta$, define the partial output
\[
 f_C(X)=\int_C m_T b_{\psi_T}(a_{\theta_T}^\top X)_+\dd\lambda_0,
 \qquad M_C=M_{\rm in},
\]
and its partial weak response
\[
 \mathsf m_{2,C}
 =\frac{\E_{P_2}[X_2(f_C)_2(X)]}{\lambda+q^2}.
\]
Let
\[
 H_2(r)=\E(Z+r)_+^2=(1+r^2)\Phi(r)+r\phi(r).
\]
This nonnegative Gaussian moment is increasing, since $H_2'=2K>0$.

\begin{proposition}[A positive weak contribution with an exact strong-tail bound]
\label{inc:AN02:arrival:v1:observables}
For the cohort in Theorem~\ref{inc:AN02:arrival:v1:main},
\begin{align}
 \mathsf m_{2,C}
 &\ge\frac{\lambda}{4(\lambda+q^2)}M_C
 \ge\frac15M_C,
 \label{inc:AN02:arrival:v1:partialresponse}\\
 \norm{f_C}_{L^2(P_1)}
 &\le M_Cq\sqrt{H_2\!\left(-\frac1{4q}\right)},
 \label{inc:AN02:arrival:v1:tailoutput}\\
 0<\frac{\E_{P_1}(f_C)_2(X)}q
 &\le M_CK\!\left(-\frac1{4q}\right).
 \label{inc:AN02:arrival:v1:partiallag}
\end{align}
Moreover $(f_C)_1(X)\le0$ and $(f_C)_2(X)\ge0$ for every $X$.
For every unit probe $\xi=(\sin\omega,-\cos\omega)$ with
$0\le\omega\le\pi/2$, $f_C(\xi)=0$: this cohort is strictly
probe-inactive at the stated time.
\end{proposition}
\begin{proof}
The Gaussian feature identity, for $a=a_{\theta_T}$, gives
\[
 \E_{P_2}[X_2(a^\top X)_+]
 =\rho qK(\rho a_2/q)+q^2a_2\Phi(\rho a_2/q)
 \ge\lambda a_2.
\]
The inequality uses $a_2>0$ and $K(r)\ge r$.
Multiplying by $b_2\ge1/2$ and using $a_2\ge1/2$ proves the first
bound in \eqref{inc:AN02:arrival:v1:partialresponse}. On the parameter
range, $\lambda\ge9/25$ and $q^2\le1/400$, so its coefficient is at
least $36/145>1/5$.

On cluster 1, $a^\top X$ is Gaussian with mean $a_1\le-1/4$ and
variance $q^2$. Hence
\[
 \norm{(a^\top X)_+}_{L^2(P_1)}
 \le q\sqrt{H_2(-1/(4q))}.
\]
Minkowski's inequality and $|b|=1$ give
\eqref{inc:AN02:arrival:v1:tailoutput}. For its first moment,
\[
 \frac{\E_{P_1}(f_C)_2(X)}q
 =\int_C m_T b_2K(a_1/q)\dd\lambda_0.
\]
All factors are positive on a positive-mass set; monotonicity of $K$
and $b_2\le1$ prove \eqref{inc:AN02:arrival:v1:partiallag}.
The coordinate signs follow from $b_1<0<b_2$ and nonnegative ReLU
features. Finally
$a^\top\xi=a_1\sin\omega-a_2\cos\omega<0$ on the specified closed
probe quadrant, so every cohort feature vanishes there.
\end{proof}

\paragraph{Scope of these signs.}
$\mathsf m_{2,C}$ is an additive contribution, not a lower bound on the
full $\mathsf m_2$: other labels can contribute with either sign.
The positive quantity in \eqref{inc:AN02:arrival:v1:partiallag} is only
this cohort's contribution to the exact mean lag, not positivity of the
full lag or of any source-resolved update. The bound in
\eqref{inc:AN02:arrival:v1:tailoutput} keeps the Gaussian tail nonzero.
Probe inactivity holds at $T_\delta$, not automatically for later times.

\section{The next missing inequality, and limits of this increment}
\label{inc:AN02:arrival:v1:next}

\paragraph{A concrete capture interface, not a claimed theorem.}
To turn this reservoir into AN3 data, the next argument must produce a
later, analytically defined $T_{\rm ent}$ and a subcohort
$C_{\rm cap}\subset C_{T_\delta}$ such that, for quantitatively controlled
$c_{\rm cap}>0$ and $R_w<\infty$,
\begin{align}
 \int_{C_{\rm cap}}m_{T_{\rm ent}}\dd\lambda_0
 &\ge c_{\rm cap}M_{\rm in},
 \label{inc:AN02:arrival:v1:openmass}\\
 \left|\frac{\pi/2-\theta_{T_{\rm ent}}}{q}\right|
 +\left|\frac{\psi_{T_{\rm ent}}-\pi/2}{q}\right|
 &\le R_w\quad\hbox{on }C_{\rm cap}.
 \label{inc:AN02:arrival:v1:openchart}
\end{align}
At the same time it must prove the strong learned-state and shape bounds,
and control every remaining label's residual contribution. The
$c_{\rm cap}$ and $R_w$ requirements are targets, not assumptions used to
prove Theorem~\ref{inc:AN02:arrival:v1:main}.
Even proving them for a fraction of this particular cohort would be one
possible route, not a claim that every future weak label must come from it.

If later capture retains mass up to a $q$-dependent factor without an
additional adverse $S_0$ power, the incoming estimate supplies the candidate
seed exponent $1+\varepsilon_R$ (or the conservative $17/16$).
Its use in AN3 would still require a proved strong-shape bound and the
nonlinear passage estimate. In particular the sign test would be
$\beta_{\rm shape}-(1+\varepsilon_R)\alpha_{\rm sh}/\lambda>0$,
with all $q$-dependent prefactors retained; neither that capture premise
nor an original-flow $\beta_{\rm shape}$ is proved here.
The displayed $qK(-R)$ prefactor and the fixed positive-$q$ contribution
$\Phi(-1/(2q))$ to $\varepsilon_R$ must not be suppressed when choosing
joint limits.

\paragraph{Why the early comparison cannot be extended by assertion.}
At its endpoint it only controls angular errors by $O(\delta)$, not
$O(q)$; at a later time with large $\overline S$, it may supply no useful
comparison at all. The existence of an incoming cohort therefore does
not prove that the residual field creates and captures it in a weak
attractor while strong learning proceeds. No state near a resident saddle
has been assumed or reached in this proof.

\paragraph{Integration proposal.}
After independent review and explicit approval, this increment can follow
\texttt{pa:targetcomparison} in the early-escape/AN2 portion of a new
checkpoint. It does not supersede any existing result. The new ingredients
are \eqref{inc:AN02:arrival:v1:tailcurvature}, the backward Jacobian bound,
and the exact initialized reservoir estimate. Do not merge it into the
current checkpoint automatically.

\begin{thebibliography}{9}
\bibitem{F}
\emph{ReLU Population Foundations}, version 2.
User-supplied \path{RELU_POPULATION_FOUNDATIONS_v2.tex}.
Dependencies: P1--P5; in particular the Gaussian feature formula (6.7)
and global characteristic existence in P5.
\bibitem{P}
\emph{Analytical progress toward initialized population reversal},
checkpoint version 1.0.
User-supplied \path{THEOREM_A_ANALYTICAL_PROGRESS.tex}.
Dependencies: \texttt{pa:setup}, \texttt{pa:earlybounds},
\texttt{pa:targetcomparison}; the AN2/AN3 contracts determine scope.
\bibitem{H}
\emph{Analytical-only theory handoff: initialized ReLU population reversal}.
User-supplied \path{THEORY_HANDOFF_ANALYTICAL_A_B.md}.
Sections 0, 2, and AN2--AN3: analytical-only proof and separate-review
workflow; prescribed initialization and entry/passage obligations.
\end{thebibliography}
\end{document}
